% ---------------------------------------------------------------------
%  6. Discusión — C (H1), A (H2, H4), B (H3), D (síntesis por pregunta)
% ---------------------------------------------------------------------
\section{Discusión: teoría frente a experimento}
\label{sec:discusion}

Esta sección contrasta cada hipótesis de la \cref{sec:hipotesis} con la evidencia de la
\cref{sec:resultados} y, a continuación, sintetiza la respuesta a cada pregunta de
investigación.

\subsection{H1: escalamiento lineal}
\label{sec:disc-h1}
\begin{aporte}{C}
\paragraph{Predicción.} Según la \cref{tab:comparacion-teorica}, los cuatro procedimientos son
lineales en $n$: la codificación de Huffman es $\Theta(n + \sigma\log\sigma)$, su
decodificación $\Theta(n\Lbar)$ con $\Lbar$ acotado, y ambas operaciones de LZW son
$\Theta(n)$. En escala log--log, la pendiente esperada es 1.

\paragraph{Desviaciones previsibles.} El análisis asintótico ignora costos fijos que pesan en
entradas pequeñas: Huffman construye el código y escribe una cabecera de 262 bytes en cada
llamada, y LZW reserva e inicializa en cada llamada su tabla hash de $2^{17}$ posiciones (dos
arreglos de enteros) o, al decodificar, cuatro arreglos de $2^{16}$ entradas. Estos costos
constantes aplanan la curva en los primeros tamaños y tienden a producir pendientes menores
que 1 si se incluyen en el ajuste; por eso el ajuste usa $n \ge 2^{16}$. En el extremo opuesto,
cuando los datos y las tablas dejan de caber en la memoria caché, el costo por byte puede
aumentar ligeramente sin que cambie la complejidad.

\paragraph{Evidencia.} \Cref{fig:e1-tiempo,fig:e1-nspb} y \cref{tab:pendientes}.
% TODO-RESULTS: contrastar H1 con las pendientes y R^2 medidos (tab:pendientes), indicar si se acepta, y explicar cualquier pendiente alejada de 1 con los costos fijos o efectos de caché descritos.
\todoresults{Contrastar H1 con las pendientes y $R^2$ de la \cref{tab:pendientes}; indicar si
se acepta y explicar cualquier pendiente alejada de 1 con los costos fijos o los efectos de
caché descritos.}
\end{aporte}

\subsection{H2: Huffman frente a la entropía}
\label{sec:disc-h2}
\begin{aporte}{A}
\paragraph{Predicción.} La tasa medida incluye la cabecera y el relleno del último byte, por lo
que
\begin{equation}
  \bpb_{\text{Huffman}} = \Lbar + \frac{8 \cdot 262}{n} + \frac{\delta}{n}, \qquad 0 \le \delta < 8,
  \label{eq:bpb-huffman}
\end{equation}
y, por la \cref{eq:cota-huffman}, tiende a un valor en $[H_0, H_0 + 1)$ cuando $n$ crece. En
E2 con $k$ potencia de dos y $n = 2^{23}$, las frecuencias empíricas son casi iguales, el árbol
resultante es completo y equilibrado, y cada símbolo recibe $\log_2 k$ bits: la teoría predice
una tasa prácticamente igual a $H_0$, con un sobrecosto de cabecera de $8\cdot 262/2^{23}$ bpb,
del orden de $2.5\times10^{-4}$. En la fuente geométrica, en cambio, las probabilidades no son
potencias de $1/2$ y cabe esperar una distancia positiva, aunque menor que un bit, respecto de
$H_0$. Sobre texto, Huffman no puede bajar de $H_0$ porque no modela el contexto.

\paragraph{Evidencia.} \Cref{fig:e2-bpb,fig:e1-bpb} y la columna $H_0$ de la \cref{tab:corpus}.
% TODO-RESULTS: verificar que todos los bpb de Huffman (E1, E2, E3) estén en [H0, H0+1) salvo el efecto de cabecera en n pequeño; cuantificar la distancia a H0 en E2 y en skewed.
\todoresults{Verificar que los bits por byte de Huffman en E1, E2 y E3 queden en
$[H_0, H_0+1)$, salvo el efecto de la cabecera en $n$ pequeño, y cuantificar la distancia a
$H_0$ en E2 y en \texttt{skewed}.}
\end{aporte}

\subsection{H3: tipo de redundancia y ganador}
\label{sec:disc-h3}
\begin{aporte}{B}
\paragraph{Predicción.} LZW solo ahorra espacio cuando sus frases son largas: cada código
cuesta entre 9 y $b$ bits, así que una frase de longitud media $m$ cuesta del orden de
$b/m$ bits por byte. En datos con palabras y frases recurrentes, el diccionario captura
secuencias largas y la tasa puede quedar por debajo de $H_0$, algo imposible para un código de
orden cero. En datos i.i.d. las frases crecen muy lentamente, porque cada extensión requiere
que un par (frase, byte) ya haya aparecido; con bytes uniformes, la mayoría de las frases
tienen uno o dos bytes y cada código de 9 a 16 bits representa poco más de un byte, de modo
que se espera expansión ($\bpb > 8$). Huffman, por el contrario, nunca supera
$\Lbar < H_0 + 1$ más la cabecera.

\paragraph{Evidencia.} \Cref{fig:e1-bpb,fig:e2-bpb} y \cref{tab:corpus}.
% TODO-RESULTS: contrastar H3 con E1 (synthtext, dickens frente a uniform, skewed), E2 y E3 (\resLzwWins de \resCorpusFiles); identificar archivos del corpus que contradigan la hipótesis y explicar por qué.
\todoresults{Contrastar H3 con E1 (\texttt{synthtext} y \texttt{dickens} frente a
\texttt{uniform} y \texttt{skewed}), con E2 y con E3 (LZW gana en \resLzwWins{} de
\resCorpusFiles{} archivos); identificar los archivos que contradigan la hipótesis y explicar
por qué.}
\end{aporte}

\subsection{H4: sensibilidad del tiempo a la entropía}
\label{sec:disc-h4}
\begin{aporte}{A}
\paragraph{Predicción.} El decodificador canónico ejecuta una iteración por bit leído, $n\Lbar$
en total; con $\Lbar \approx \log_2 k$ en E2, su tiempo debería crecer aproximadamente de forma
lineal con $\log_2 k$. El codificador de Huffman escribe un campo por símbolo con
\texttt{BitWriter}, cuyo costo depende solo débilmente de $\ell_s$, por lo que se espera
menor sensibilidad. La decodificación de LZW escribe cada byte de salida una vez; la entropía
solo afecta al número de códigos, que crece cuando las frases se acortan, de modo que se espera
una dependencia mucho más débil que la de Huffman.

\paragraph{Evidencia.} \Cref{fig:e2-tiempo}.
% TODO-RESULTS: contrastar H4 con E2 (pendiente aproximada del tiempo de decodificación de Huffman frente a log2 k; comparar con LZW usando \resEntropyHuffDecMsKTwo, \resEntropyHuffDecMsKTwoFiftySix, \resEntropyLzwDecMsKTwo, \resEntropyLzwDecMsKTwoFiftySix).
\todoresults{Contrastar H4 con la \cref{fig:e2-tiempo}: tendencia del tiempo de decodificación
de Huffman frente a $\log_2 k$ comparada con la de LZW.}
\end{aporte}

\subsection{Síntesis por pregunta de investigación}
\label{sec:disc-rq}
\begin{aporte}{D}
\paragraph{RQ1 (escalamiento).} Se responde con H1 y H4: la teoría predice crecimiento lineal
para ambos algoritmos, con una constante que en la decodificación de Huffman depende de la
longitud media del código.
% TODO-RESULTS: respuesta a RQ1 basada en H1 y H4.
\todoresults{Responder RQ1 a partir de la discusión de H1 y H4.}

\paragraph{RQ2 (compresión, entropía y tipo de dato).} Se responde con H2 y H3: Huffman está
acotado por $H_0$ y no puede aprovechar repeticiones, mientras que LZW puede bajar de $H_0$ en
datos con frases repetidas, pero pierde ---y puede expandir--- en fuentes sin memoria.
% TODO-RESULTS: respuesta a RQ2 basada en H2, H3 y E3.
\todoresults{Responder RQ2 a partir de H2, H3 y del agregado del corpus.}

\paragraph{RQ3 (tamaño del diccionario).} Un diccionario mayor retiene más frases y reinicia con
menos frecuencia, pero cada código cuesta más bits; la teoría no fija un óptimo universal, que
depende de cuán repetitivos y estacionarios sean los datos~\cite{cheng2021}.
% TODO-RESULTS: respuesta a RQ3 basada en E4 (\resLzwBitsBestDickens, \resLzwBitsBestXml, \resLzwBitsBestEcoli).
\todoresults{Responder RQ3 a partir de E4 (mejor ancho: \texttt{dickens}
$b = \resLzwBitsBestDickens$, \texttt{xml} $b = \resLzwBitsBestXml$, \texttt{E.coli}
$b = \resLzwBitsBestEcoli$).}
\end{aporte}
